\textbf{A continuación se presenta la descomposición departamental del crecimiento de la productividad laboral.} Para cada departamento se reportan el PIB real, el número de ocupados, las horas semanales por trabajador, el PIB por trabajador y el PIB por hora trabajada al inicio y al final del periodo disponible.

\subsection{Quindío}
\textbf{Entre 2014 y 2024, el PIB por hora trabajada de Quindío creció 8,21\% anual.} El PIB por trabajador creció 7,21\% anual, mientras que las horas semanales por trabajador cambiaron -0,92\% anual. Esta diferencia muestra si la productividad medida por trabajador se mueve en línea con la productividad por hora o si está afectada por cambios en la intensidad horaria.
\begin{table}[H]
\centering
\caption{Quindío: PIB, ocupados, horas y productividad laboral, 2014--2024pr}
\label{tab:departamento_quindio_productividad}
\scriptsize
\begin{tabular}{lrrr}
\toprule
Indicador & 2014 & 2024pr & Crec. anual \\
\midrule
PIB real & 6,1 & 8,0 & 2,70\% \\
Ocupados & 0,36 & 0,24 & -4,21\% \\
Horas semanales por trabajador & 47,3 & 43,1 & -0,92\% \\
PIB por trabajador & 16,8 & 33,8 & 7,21\% \\
PIB por hora & 6,8 & 15,1 & 8,21\% \\
\bottomrule
\end{tabular}
\end{table}

\subsection{Putumayo}
\textbf{Entre 2014 y 2024, el PIB por hora trabajada de Putumayo creció 6,62\% anual.} El PIB por trabajador creció 6,27\% anual, mientras que las horas semanales por trabajador cambiaron -0,33\% anual. Esta diferencia muestra si la productividad medida por trabajador se mueve en línea con la productividad por hora o si está afectada por cambios en la intensidad horaria.
\begin{table}[H]
\centering
\caption{Putumayo: PIB, ocupados, horas y productividad laboral, 2014--2024pr}
\label{tab:departamento_putumayo_productividad}
\scriptsize
\begin{tabular}{lrrr}
\toprule
Indicador & 2014 & 2024pr & Crec. anual \\
\midrule
PIB real & 3,8 & 3,4 & -0,92\% \\
Ocupados & 0,02 & 0,01 & -6,77\% \\
Horas semanales por trabajador & 43,6 & 42,1 & -0,33\% \\
PIB por trabajador & 241,1 & 443,0 & 6,27\% \\
PIB por hora & 106,4 & 202,1 & 6,62\% \\
\bottomrule
\end{tabular}
\end{table}

\subsection{Risaralda}
\textbf{Entre 2014 y 2024, el PIB por hora trabajada de Risaralda creció 4,14\% anual.} El PIB por trabajador creció 3,37\% anual, mientras que las horas semanales por trabajador cambiaron -0,74\% anual. Esta diferencia muestra si la productividad medida por trabajador se mueve en línea con la productividad por hora o si está afectada por cambios en la intensidad horaria.
\begin{table}[H]
\centering
\caption{Risaralda: PIB, ocupados, horas y productividad laboral, 2014--2024pr}
\label{tab:departamento_risaralda_productividad}
\scriptsize
\begin{tabular}{lrrr}
\toprule
Indicador & 2014 & 2024pr & Crec. anual \\
\midrule
PIB real & 12,4 & 16,3 & 2,83\% \\
Ocupados & 0,45 & 0,42 & -0,52\% \\
Horas semanales por trabajador & 48,9 & 45,4 & -0,74\% \\
PIB por trabajador & 27,8 & 38,7 & 3,37\% \\
PIB por hora & 10,9 & 16,4 & 4,14\% \\
\bottomrule
\end{tabular}
\end{table}

\subsection{Caquetá}
\textbf{Entre 2014 y 2024, el PIB por hora trabajada de Caquetá creció 3,88\% anual.} El PIB por trabajador creció 3,51\% anual, mientras que las horas semanales por trabajador cambiaron -0,35\% anual. Esta diferencia muestra si la productividad medida por trabajador se mueve en línea con la productividad por hora o si está afectada por cambios en la intensidad horaria.
\begin{table}[H]
\centering
\caption{Caquetá: PIB, ocupados, horas y productividad laboral, 2014--2024pr}
\label{tab:departamento_caqueta_productividad}
\scriptsize
\begin{tabular}{lrrr}
\toprule
Indicador & 2014 & 2024pr & Crec. anual \\
\midrule
PIB real & 3,3 & 3,9 & 1,60\% \\
Ocupados & 0,19 & 0,16 & -1,85\% \\
Horas semanales por trabajador & 47,3 & 45,7 & -0,35\% \\
PIB por trabajador & 17,6 & 24,9 & 3,51\% \\
PIB por hora & 7,2 & 10,5 & 3,88\% \\
\bottomrule
\end{tabular}
\end{table}

\subsection{Caldas}
\textbf{Entre 2014 y 2024, el PIB por hora trabajada de Caldas creció 2,81\% anual.} El PIB por trabajador creció 2,36\% anual, mientras que las horas semanales por trabajador cambiaron -0,44\% anual. Esta diferencia muestra si la productividad medida por trabajador se mueve en línea con la productividad por hora o si está afectada por cambios en la intensidad horaria.
\begin{table}[H]
\centering
\caption{Caldas: PIB, ocupados, horas y productividad laboral, 2014--2024pr}
\label{tab:departamento_caldas_productividad}
\scriptsize
\begin{tabular}{lrrr}
\toprule
Indicador & 2014 & 2024pr & Crec. anual \\
\midrule
PIB real & 12,1 & 15,7 & 2,66\% \\
Ocupados & 0,42 & 0,43 & 0,30\% \\
Horas semanales por trabajador & 47,3 & 45,2 & -0,44\% \\
PIB por trabajador & 29,0 & 36,6 & 2,36\% \\
PIB por hora & 11,8 & 15,6 & 2,81\% \\
\bottomrule
\end{tabular}
\end{table}

\subsection{Sucre}
\textbf{Entre 2014 y 2024, el PIB por hora trabajada de Sucre creció 2,79\% anual.} El PIB por trabajador creció 2,03\% anual, mientras que las horas semanales por trabajador cambiaron -0,74\% anual. Esta diferencia muestra si la productividad medida por trabajador se mueve en línea con la productividad por hora o si está afectada por cambios en la intensidad horaria.
\begin{table}[H]
\centering
\caption{Sucre: PIB, ocupados, horas y productividad laboral, 2014--2024pr}
\label{tab:departamento_sucre_productividad}
\scriptsize
\begin{tabular}{lrrr}
\toprule
Indicador & 2014 & 2024pr & Crec. anual \\
\midrule
PIB real & 6,5 & 8,1 & 2,30\% \\
Ocupados & 0,36 & 0,37 & 0,27\% \\
Horas semanales por trabajador & 43,9 & 40,8 & -0,74\% \\
PIB por trabajador & 18,0 & 22,0 & 2,03\% \\
PIB por hora & 7,9 & 10,4 & 2,79\% \\
\bottomrule
\end{tabular}
\end{table}

\subsection{San Andrés y Providencia}
\textbf{Entre 2014 y 2024, el PIB por hora trabajada de San Andrés y Providencia creció 2,78\% anual.} El PIB por trabajador creció 2,86\% anual, mientras que las horas semanales por trabajador cambiaron 0,08\% anual. Esta diferencia muestra si la productividad medida por trabajador se mueve en línea con la productividad por hora o si está afectada por cambios en la intensidad horaria.
\begin{table}[H]
\centering
\caption{San Andrés y Providencia: PIB, ocupados, horas y productividad laboral, 2014--2024pr}
\label{tab:departamento_san_andres_y_providencia_productividad}
\scriptsize
\begin{tabular}{lrrr}
\toprule
Indicador & 2014 & 2024pr & Crec. anual \\
\midrule
PIB real & 1,2 & 1,5 & 2,48\% \\
Ocupados & 0,02 & 0,02 & -0,37\% \\
Horas semanales por trabajador & 47,2 & 47,6 & 0,08\% \\
PIB por trabajador & 62,1 & 82,4 & 2,86\% \\
PIB por hora & 25,3 & 33,3 & 2,78\% \\
\bottomrule
\end{tabular}
\end{table}

\subsection{Bogotá D.C.}
\textbf{Entre 2014 y 2024, el PIB por hora trabajada de Bogotá D.C. creció 2,44\% anual.} El PIB por trabajador creció 2,05\% anual, mientras que las horas semanales por trabajador cambiaron -0,38\% anual. Esta diferencia muestra si la productividad medida por trabajador se mueve en línea con la productividad por hora o si está afectada por cambios en la intensidad horaria.
\begin{table}[H]
\centering
\caption{Bogotá D.C.: PIB, ocupados, horas y productividad laboral, 2014--2024pr}
\label{tab:departamento_bogota_dc_productividad}
\scriptsize
\begin{tabular}{lrrr}
\toprule
Indicador & 2014 & 2024pr & Crec. anual \\
\midrule
PIB real & 199,0 & 267,1 & 2,99\% \\
Ocupados & 3,62 & 3,96 & 0,91\% \\
Horas semanales por trabajador & 47,1 & 45,3 & -0,38\% \\
PIB por trabajador & 55,0 & 67,4 & 2,05\% \\
PIB por hora & 22,5 & 28,6 & 2,44\% \\
\bottomrule
\end{tabular}
\end{table}

\subsection{Boyacá}
\textbf{Entre 2014 y 2024, el PIB por hora trabajada de Boyacá creció 2,16\% anual.} El PIB por trabajador creció 1,49\% anual, mientras que las horas semanales por trabajador cambiaron -0,66\% anual. Esta diferencia muestra si la productividad medida por trabajador se mueve en línea con la productividad por hora o si está afectada por cambios en la intensidad horaria.
\begin{table}[H]
\centering
\caption{Boyacá: PIB, ocupados, horas y productividad laboral, 2014--2024pr}
\label{tab:departamento_boyaca_productividad}
\scriptsize
\begin{tabular}{lrrr}
\toprule
Indicador & 2014 & 2024pr & Crec. anual \\
\midrule
PIB real & 21,0 & 25,2 & 1,84\% \\
Ocupados & 0,51 & 0,53 & 0,35\% \\
Horas semanales por trabajador & 43,2 & 40,4 & -0,66\% \\
PIB por trabajador & 41,0 & 47,5 & 1,49\% \\
PIB por hora & 18,3 & 22,6 & 2,16\% \\
\bottomrule
\end{tabular}
\end{table}

\subsection{Tolima}
\textbf{Entre 2014 y 2024, el PIB por hora trabajada de Tolima creció 1,75\% anual.} El PIB por trabajador creció 1,76\% anual, mientras que las horas semanales por trabajador cambiaron 0,01\% anual. Esta diferencia muestra si la productividad medida por trabajador se mueve en línea con la productividad por hora o si está afectada por cambios en la intensidad horaria.
\begin{table}[H]
\centering
\caption{Tolima: PIB, ocupados, horas y productividad laboral, 2014--2024pr}
\label{tab:departamento_tolima_productividad}
\scriptsize
\begin{tabular}{lrrr}
\toprule
Indicador & 2014 & 2024pr & Crec. anual \\
\midrule
PIB real & 17,2 & 20,3 & 1,69\% \\
Ocupados & 0,51 & 0,51 & -0,07\% \\
Horas semanales por trabajador & 45,4 & 45,5 & 0,01\% \\
PIB por trabajador & 33,6 & 40,0 & 1,76\% \\
PIB por hora & 14,2 & 16,9 & 1,75\% \\
\bottomrule
\end{tabular}
\end{table}

\subsection{Valle del Cauca}
\textbf{Entre 2014 y 2024, el PIB por hora trabajada de Valle del Cauca creció 1,72\% anual.} El PIB por trabajador creció 1,39\% anual, mientras que las horas semanales por trabajador cambiaron -0,33\% anual. Esta diferencia muestra si la productividad medida por trabajador se mueve en línea con la productividad por hora o si está afectada por cambios en la intensidad horaria.
\begin{table}[H]
\centering
\caption{Valle del Cauca: PIB, ocupados, horas y productividad laboral, 2014--2024pr}
\label{tab:departamento_valle_del_cauca_productividad}
\scriptsize
\begin{tabular}{lrrr}
\toprule
Indicador & 2014 & 2024pr & Crec. anual \\
\midrule
PIB real & 75,9 & 99,5 & 2,75\% \\
Ocupados & 1,60 & 1,83 & 1,34\% \\
Horas semanales por trabajador & 46,2 & 44,7 & -0,33\% \\
PIB por trabajador & 47,5 & 54,5 & 1,39\% \\
PIB por hora & 19,8 & 23,5 & 1,72\% \\
\bottomrule
\end{tabular}
\end{table}

\subsection{Santander}
\textbf{Entre 2014 y 2024, el PIB por hora trabajada de Santander creció 1,68\% anual.} El PIB por trabajador creció 1,40\% anual, mientras que las horas semanales por trabajador cambiaron -0,28\% anual. Esta diferencia muestra si la productividad medida por trabajador se mueve en línea con la productividad por hora o si está afectada por cambios en la intensidad horaria.
\begin{table}[H]
\centering
\caption{Santander: PIB, ocupados, horas y productividad laboral, 2014--2024pr}
\label{tab:departamento_santander_productividad}
\scriptsize
\begin{tabular}{lrrr}
\toprule
Indicador & 2014 & 2024pr & Crec. anual \\
\midrule
PIB real & 52,1 & 62,6 & 1,84\% \\
Ocupados & 0,97 & 1,01 & 0,43\% \\
Horas semanales por trabajador & 46,9 & 45,6 & -0,28\% \\
PIB por trabajador & 53,9 & 62,0 & 1,40\% \\
PIB por hora & 22,1 & 26,1 & 1,68\% \\
\bottomrule
\end{tabular}
\end{table}

\subsection{Chocó}
\textbf{Entre 2014 y 2024, el PIB por hora trabajada de Chocó creció 0,89\% anual.} El PIB por trabajador creció 1,96\% anual, mientras que las horas semanales por trabajador cambiaron 1,06\% anual. Esta diferencia muestra si la productividad medida por trabajador se mueve en línea con la productividad por hora o si está afectada por cambios en la intensidad horaria.
\begin{table}[H]
\centering
\caption{Chocó: PIB, ocupados, horas y productividad laboral, 2014--2024pr}
\label{tab:departamento_choco_productividad}
\scriptsize
\begin{tabular}{lrrr}
\toprule
Indicador & 2014 & 2024pr & Crec. anual \\
\midrule
PIB real & 3,5 & 3,8 & 0,83\% \\
Ocupados & 0,15 & 0,13 & -1,11\% \\
Horas semanales por trabajador & 41,5 & 46,1 & 1,06\% \\
PIB por trabajador & 23,6 & 28,6 & 1,96\% \\
PIB por hora & 10,9 & 11,9 & 0,89\% \\
\bottomrule
\end{tabular}
\end{table}

\subsection{Bolívar}
\textbf{Entre 2014 y 2024, el PIB por hora trabajada de Bolívar creció 0,82\% anual.} El PIB por trabajador creció 0,61\% anual, mientras que las horas semanales por trabajador cambiaron -0,21\% anual. Esta diferencia muestra si la productividad medida por trabajador se mueve en línea con la productividad por hora o si está afectada por cambios en la intensidad horaria.
\begin{table}[H]
\centering
\caption{Bolívar: PIB, ocupados, horas y productividad laboral, 2014--2024pr}
\label{tab:departamento_bolivar_productividad}
\scriptsize
\begin{tabular}{lrrr}
\toprule
Indicador & 2014 & 2024pr & Crec. anual \\
\midrule
PIB real & 27,8 & 35,9 & 2,59\% \\
Ocupados & 0,66 & 0,80 & 1,97\% \\
Horas semanales por trabajador & 44,8 & 43,9 & -0,21\% \\
PIB por trabajador & 42,4 & 45,1 & 0,61\% \\
PIB por hora & 18,2 & 19,8 & 0,82\% \\
\bottomrule
\end{tabular}
\end{table}

\subsection{Norte de Santander}
\textbf{Entre 2014 y 2024, el PIB por hora trabajada de Norte de Santander creció 0,79\% anual.} El PIB por trabajador creció 0,94\% anual, mientras que las horas semanales por trabajador cambiaron 0,15\% anual. Esta diferencia muestra si la productividad medida por trabajador se mueve en línea con la productividad por hora o si está afectada por cambios en la intensidad horaria.
\begin{table}[H]
\centering
\caption{Norte de Santander: PIB, ocupados, horas y productividad laboral, 2014--2024pr}
\label{tab:departamento_norte_de_santander_productividad}
\scriptsize
\begin{tabular}{lrrr}
\toprule
Indicador & 2014 & 2024pr & Crec. anual \\
\midrule
PIB real & 12,3 & 15,6 & 2,46\% \\
Ocupados & 0,55 & 0,64 & 1,51\% \\
Horas semanales por trabajador & 46,0 & 46,7 & 0,15\% \\
PIB por trabajador & 22,2 & 24,3 & 0,94\% \\
PIB por hora & 9,3 & 10,0 & 0,79\% \\
\bottomrule
\end{tabular}
\end{table}

\subsection{Vichada}
\textbf{Entre 2014 y 2024, el PIB por hora trabajada de Vichada creció 0,77\% anual.} El PIB por trabajador creció 1,33\% anual, mientras que las horas semanales por trabajador cambiaron 0,56\% anual. Esta diferencia muestra si la productividad medida por trabajador se mueve en línea con la productividad por hora o si está afectada por cambios en la intensidad horaria.
\begin{table}[H]
\centering
\caption{Vichada: PIB, ocupados, horas y productividad laboral, 2014--2024pr}
\label{tab:departamento_vichada_productividad}
\scriptsize
\begin{tabular}{lrrr}
\toprule
Indicador & 2014 & 2024pr & Crec. anual \\
\midrule
PIB real & 0,5 & 0,6 & 2,71\% \\
Ocupados & 0,01 & 0,01 & 1,36\% \\
Horas semanales por trabajador & 44,0 & 46,5 & 0,56\% \\
PIB por trabajador & 80,8 & 92,2 & 1,33\% \\
PIB por hora & 35,3 & 38,1 & 0,77\% \\
\bottomrule
\end{tabular}
\end{table}

\subsection{Nariño}
\textbf{Entre 2014 y 2024, el PIB por hora trabajada de Nariño creció 0,46\% anual.} El PIB por trabajador creció -0,49\% anual, mientras que las horas semanales por trabajador cambiaron -0,94\% anual. Esta diferencia muestra si la productividad medida por trabajador se mueve en línea con la productividad por hora o si está afectada por cambios en la intensidad horaria.
\begin{table}[H]
\centering
\caption{Nariño: PIB, ocupados, horas y productividad laboral, 2014--2024pr}
\label{tab:departamento_narino_productividad}
\scriptsize
\begin{tabular}{lrrr}
\toprule
Indicador & 2014 & 2024pr & Crec. anual \\
\midrule
PIB real & 11,6 & 14,7 & 2,40\% \\
Ocupados & 0,62 & 0,82 & 2,90\% \\
Horas semanales por trabajador & 41,2 & 37,5 & -0,94\% \\
PIB por trabajador & 18,7 & 17,8 & -0,49\% \\
PIB por hora & 8,7 & 9,1 & 0,46\% \\
\bottomrule
\end{tabular}
\end{table}

\subsection{Cauca}
\textbf{Entre 2014 y 2024, el PIB por hora trabajada de Cauca creció 0,32\% anual.} El PIB por trabajador creció -0,21\% anual, mientras que las horas semanales por trabajador cambiaron -0,53\% anual. Esta diferencia muestra si la productividad medida por trabajador se mueve en línea con la productividad por hora o si está afectada por cambios en la intensidad horaria.
\begin{table}[H]
\centering
\caption{Cauca: PIB, ocupados, horas y productividad laboral, 2014--2024pr}
\label{tab:departamento_cauca_productividad}
\scriptsize
\begin{tabular}{lrrr}
\toprule
Indicador & 2014 & 2024pr & Crec. anual \\
\midrule
PIB real & 13,9 & 17,4 & 2,28\% \\
Ocupados & 0,51 & 0,65 & 2,50\% \\
Horas semanales por trabajador & 39,4 & 37,4 & -0,53\% \\
PIB por trabajador & 27,2 & 26,6 & -0,21\% \\
PIB por hora & 13,3 & 13,7 & 0,32\% \\
\bottomrule
\end{tabular}
\end{table}

\subsection{Córdoba}
\textbf{Entre 2014 y 2024, el PIB por hora trabajada de Córdoba creció 0,25\% anual.} El PIB por trabajador creció 0,16\% anual, mientras que las horas semanales por trabajador cambiaron -0,09\% anual. Esta diferencia muestra si la productividad medida por trabajador se mueve en línea con la productividad por hora o si está afectada por cambios en la intensidad horaria.
\begin{table}[H]
\centering
\caption{Córdoba: PIB, ocupados, horas y productividad laboral, 2014--2024pr}
\label{tab:departamento_cordoba_productividad}
\scriptsize
\begin{tabular}{lrrr}
\toprule
Indicador & 2014 & 2024pr & Crec. anual \\
\midrule
PIB real & 13,7 & 17,0 & 2,21\% \\
Ocupados & 0,63 & 0,77 & 2,05\% \\
Horas semanales por trabajador & 40,1 & 39,7 & -0,09\% \\
PIB por trabajador & 21,7 & 22,0 & 0,16\% \\
PIB por hora & 10,4 & 10,7 & 0,25\% \\
\bottomrule
\end{tabular}
\end{table}

\subsection{Antioquia}
\textbf{Entre 2014 y 2024, el PIB por hora trabajada de Antioquia creció 0,21\% anual.} El PIB por trabajador creció -0,09\% anual, mientras que las horas semanales por trabajador cambiaron -0,30\% anual. Esta diferencia muestra si la productividad medida por trabajador se mueve en línea con la productividad por hora o si está afectada por cambios en la intensidad horaria.
\begin{table}[H]
\centering
\caption{Antioquia: PIB, ocupados, horas y productividad laboral, 2014--2024pr}
\label{tab:departamento_antioquia_productividad}
\scriptsize
\begin{tabular}{lrrr}
\toprule
Indicador & 2014 & 2024pr & Crec. anual \\
\midrule
PIB real & 111,1 & 149,2 & 2,99\% \\
Ocupados & 2,31 & 3,13 & 3,08\% \\
Horas semanales por trabajador & 46,4 & 45,1 & -0,30\% \\
PIB por trabajador & 48,1 & 47,7 & -0,09\% \\
PIB por hora & 19,9 & 20,3 & 0,21\% \\
\bottomrule
\end{tabular}
\end{table}

\subsection{Amazonas}
\textbf{Entre 2014 y 2024, el PIB por hora trabajada de Amazonas creció 0,20\% anual.} El PIB por trabajador creció -0,09\% anual, mientras que las horas semanales por trabajador cambiaron -0,30\% anual. Esta diferencia muestra si la productividad medida por trabajador se mueve en línea con la productividad por hora o si está afectada por cambios en la intensidad horaria.
\begin{table}[H]
\centering
\caption{Amazonas: PIB, ocupados, horas y productividad laboral, 2014--2024pr}
\label{tab:departamento_amazonas_productividad}
\scriptsize
\begin{tabular}{lrrr}
\toprule
Indicador & 2014 & 2024pr & Crec. anual \\
\midrule
PIB real & 0,6 & 0,7 & 2,55\% \\
Ocupados & 0,01 & 0,01 & 2,64\% \\
Horas semanales por trabajador & 45,7 & 44,3 & -0,30\% \\
PIB por trabajador & 50,7 & 50,2 & -0,09\% \\
PIB por hora & 21,3 & 21,8 & 0,20\% \\
\bottomrule
\end{tabular}
\end{table}

\subsection{Magdalena}
\textbf{Entre 2014 y 2024, el PIB por hora trabajada de Magdalena creció 0,17\% anual.} El PIB por trabajador creció 0,00\% anual, mientras que las horas semanales por trabajador cambiaron -0,16\% anual. Esta diferencia muestra si la productividad medida por trabajador se mueve en línea con la productividad por hora o si está afectada por cambios en la intensidad horaria.
\begin{table}[H]
\centering
\caption{Magdalena: PIB, ocupados, horas y productividad laboral, 2014--2024pr}
\label{tab:departamento_magdalena_productividad}
\scriptsize
\begin{tabular}{lrrr}
\toprule
Indicador & 2014 & 2024pr & Crec. anual \\
\midrule
PIB real & 10,1 & 13,1 & 2,63\% \\
Ocupados & 0,42 & 0,55 & 2,62\% \\
Horas semanales por trabajador & 45,5 & 44,7 & -0,16\% \\
PIB por trabajador & 23,9 & 23,9 & 0,00\% \\
PIB por hora & 10,1 & 10,3 & 0,17\% \\
\bottomrule
\end{tabular}
\end{table}

\subsection{Guaviare}
\textbf{Entre 2014 y 2024, el PIB por hora trabajada de Guaviare creció -0,14\% anual.} El PIB por trabajador creció -1,00\% anual, mientras que las horas semanales por trabajador cambiaron -0,86\% anual. Esta diferencia muestra si la productividad medida por trabajador se mueve en línea con la productividad por hora o si está afectada por cambios en la intensidad horaria.
\begin{table}[H]
\centering
\caption{Guaviare: PIB, ocupados, horas y productividad laboral, 2014--2024pr}
\label{tab:departamento_guaviare_productividad}
\scriptsize
\begin{tabular}{lrrr}
\toprule
Indicador & 2014 & 2024pr & Crec. anual \\
\midrule
PIB real & 0,6 & 0,8 & 2,28\% \\
Ocupados & 0,01 & 0,02 & 3,31\% \\
Horas semanales por trabajador & 48,1 & 44,1 & -0,86\% \\
PIB por trabajador & 43,2 & 39,1 & -1,00\% \\
PIB por hora & 17,3 & 17,0 & -0,14\% \\
\bottomrule
\end{tabular}
\end{table}

\subsection{La Guajira}
\textbf{Entre 2014 y 2024, el PIB por hora trabajada de La Guajira creció -0,39\% anual.} El PIB por trabajador creció -0,71\% anual, mientras que las horas semanales por trabajador cambiaron -0,32\% anual. Esta diferencia muestra si la productividad medida por trabajador se mueve en línea con la productividad por hora o si está afectada por cambios en la intensidad horaria.
\begin{table}[H]
\centering
\caption{La Guajira: PIB, ocupados, horas y productividad laboral, 2014--2024pr}
\label{tab:departamento_la_guajira_productividad}
\scriptsize
\begin{tabular}{lrrr}
\toprule
Indicador & 2014 & 2024pr & Crec. anual \\
\midrule
PIB real & 8,7 & 8,2 & -0,59\% \\
Ocupados & 0,33 & 0,33 & 0,12\% \\
Horas semanales por trabajador & 44,4 & 43,0 & -0,32\% \\
PIB por trabajador & 26,6 & 24,8 & -0,71\% \\
PIB por hora & 11,6 & 11,1 & -0,39\% \\
\bottomrule
\end{tabular}
\end{table}

\subsection{Arauca}
\textbf{Entre 2014 y 2024, el PIB por hora trabajada de Arauca creció -0,84\% anual.} El PIB por trabajador creció -0,28\% anual, mientras que las horas semanales por trabajador cambiaron 0,56\% anual. Esta diferencia muestra si la productividad medida por trabajador se mueve en línea con la productividad por hora o si está afectada por cambios en la intensidad horaria.
\begin{table}[H]
\centering
\caption{Arauca: PIB, ocupados, horas y productividad laboral, 2014--2024pr}
\label{tab:departamento_arauca_productividad}
\scriptsize
\begin{tabular}{lrrr}
\toprule
Indicador & 2014 & 2024pr & Crec. anual \\
\midrule
PIB real & 4,3 & 5,0 & 1,42\% \\
Ocupados & 0,02 & 0,03 & 1,70\% \\
Horas semanales por trabajador & 46,7 & 49,3 & 0,56\% \\
PIB por trabajador & 176,1 & 171,2 & -0,28\% \\
PIB por hora & 72,6 & 66,8 & -0,84\% \\
\bottomrule
\end{tabular}
\end{table}

\subsection{Huila}
\textbf{Entre 2014 y 2024, el PIB por hora trabajada de Huila creció -0,86\% anual.} El PIB por trabajador creció -1,31\% anual, mientras que las horas semanales por trabajador cambiaron -0,45\% anual. Esta diferencia muestra si la productividad medida por trabajador se mueve en línea con la productividad por hora o si está afectada por cambios en la intensidad horaria.
\begin{table}[H]
\centering
\caption{Huila: PIB, ocupados, horas y productividad laboral, 2014--2024pr}
\label{tab:departamento_huila_productividad}
\scriptsize
\begin{tabular}{lrrr}
\toprule
Indicador & 2014 & 2024pr & Crec. anual \\
\midrule
PIB real & 13,9 & 15,4 & 0,99\% \\
Ocupados & 0,39 & 0,49 & 2,33\% \\
Horas semanales por trabajador & 44,7 & 42,7 & -0,45\% \\
PIB por trabajador & 36,1 & 31,6 & -1,31\% \\
PIB por hora & 15,5 & 14,2 & -0,86\% \\
\bottomrule
\end{tabular}
\end{table}

\subsection{Cesar}
\textbf{Entre 2014 y 2024, el PIB por hora trabajada de Cesar creció -1,07\% anual.} El PIB por trabajador creció -1,61\% anual, mientras que las horas semanales por trabajador cambiaron -0,55\% anual. Esta diferencia muestra si la productividad medida por trabajador se mueve en línea con la productividad por hora o si está afectada por cambios en la intensidad horaria.
\begin{table}[H]
\centering
\caption{Cesar: PIB, ocupados, horas y productividad laboral, 2014--2024pr}
\label{tab:departamento_cesar_productividad}
\scriptsize
\begin{tabular}{lrrr}
\toprule
Indicador & 2014 & 2024pr & Crec. anual \\
\midrule
PIB real & 14,4 & 15,2 & 0,54\% \\
Ocupados & 0,38 & 0,47 & 2,19\% \\
Horas semanales por trabajador & 47,8 & 45,2 & -0,55\% \\
PIB por trabajador & 38,3 & 32,6 & -1,61\% \\
PIB por hora & 15,4 & 13,8 & -1,07\% \\
\bottomrule
\end{tabular}
\end{table}

\subsection{Guainía}
\textbf{Entre 2014 y 2024, el PIB por hora trabajada de Guainía creció -1,23\% anual.} El PIB por trabajador creció -1,46\% anual, mientras que las horas semanales por trabajador cambiaron -0,23\% anual. Esta diferencia muestra si la productividad medida por trabajador se mueve en línea con la productividad por hora o si está afectada por cambios en la intensidad horaria.
\begin{table}[H]
\centering
\caption{Guainía: PIB, ocupados, horas y productividad laboral, 2014--2024pr}
\label{tab:departamento_guainia_productividad}
\scriptsize
\begin{tabular}{lrrr}
\toprule
Indicador & 2014 & 2024pr & Crec. anual \\
\midrule
PIB real & 0,3 & 0,3 & 1,56\% \\
Ocupados & 0,01 & 0,01 & 3,06\% \\
Horas semanales por trabajador & 44,9 & 43,9 & -0,23\% \\
PIB por trabajador & 44,9 & 38,8 & -1,46\% \\
PIB por hora & 19,2 & 17,0 & -1,23\% \\
\bottomrule
\end{tabular}
\end{table}

\subsection{Atlántico}
\textbf{Entre 2014 y 2024, el PIB por hora trabajada de Atlántico creció -1,33\% anual.} El PIB por trabajador creció -1,77\% anual, mientras que las horas semanales por trabajador cambiaron -0,44\% anual. Esta diferencia muestra si la productividad medida por trabajador se mueve en línea con la productividad por hora o si está afectada por cambios en la intensidad horaria.
\begin{table}[H]
\centering
\caption{Atlántico: PIB, ocupados, horas y productividad laboral, 2014--2024pr}
\label{tab:departamento_atlantico_productividad}
\scriptsize
\begin{tabular}{lrrr}
\toprule
Indicador & 2014 & 2024pr & Crec. anual \\
\midrule
PIB real & 33,8 & 45,1 & 2,94\% \\
Ocupados & 0,76 & 1,21 & 4,79\% \\
Horas semanales por trabajador & 47,0 & 45,0 & -0,44\% \\
PIB por trabajador & 44,6 & 37,3 & -1,77\% \\
PIB por hora & 18,3 & 16,0 & -1,33\% \\
\bottomrule
\end{tabular}
\end{table}

\subsection{Cundinamarca}
\textbf{Entre 2014 y 2024, el PIB por hora trabajada de Cundinamarca creció -1,53\% anual.} El PIB por trabajador creció -1,69\% anual, mientras que las horas semanales por trabajador cambiaron -0,15\% anual. Esta diferencia muestra si la productividad medida por trabajador se mueve en línea con la productividad por hora o si está afectada por cambios en la intensidad horaria.
\begin{table}[H]
\centering
\caption{Cundinamarca: PIB, ocupados, horas y productividad laboral, 2014--2024pr}
\label{tab:departamento_cundinamarca_productividad}
\scriptsize
\begin{tabular}{lrrr}
\toprule
Indicador & 2014 & 2024pr & Crec. anual \\
\midrule
PIB real & 46,2 & 60,5 & 2,73\% \\
Ocupados & 0,97 & 1,50 & 4,49\% \\
Horas semanales por trabajador & 45,4 & 44,7 & -0,15\% \\
PIB por trabajador & 47,9 & 40,4 & -1,69\% \\
PIB por hora & 20,3 & 17,4 & -1,53\% \\
\bottomrule
\end{tabular}
\end{table}

\subsection{Meta}
\textbf{Entre 2014 y 2024, el PIB por hora trabajada de Meta creció -2,07\% anual.} El PIB por trabajador creció -2,67\% anual, mientras que las horas semanales por trabajador cambiaron -0,61\% anual. Esta diferencia muestra si la productividad medida por trabajador se mueve en línea con la productividad por hora o si está afectada por cambios en la intensidad horaria.
\begin{table}[H]
\centering
\caption{Meta: PIB, ocupados, horas y productividad laboral, 2014--2024pr}
\label{tab:departamento_meta_productividad}
\scriptsize
\begin{tabular}{lrrr}
\toprule
Indicador & 2014 & 2024pr & Crec. anual \\
\midrule
PIB real & 29,9 & 32,1 & 0,71\% \\
Ocupados & 0,36 & 0,51 & 3,47\% \\
Horas semanales por trabajador & 49,8 & 46,8 & -0,61\% \\
PIB por trabajador & 82,7 & 63,1 & -2,67\% \\
PIB por hora & 31,9 & 25,9 & -2,07\% \\
\bottomrule
\end{tabular}
\end{table}

\subsection{Vaupés}
\textbf{Entre 2014 y 2024, el PIB por hora trabajada de Vaupés creció -2,75\% anual.} El PIB por trabajador creció -3,31\% anual, mientras que las horas semanales por trabajador cambiaron -0,57\% anual. Esta diferencia muestra si la productividad medida por trabajador se mueve en línea con la productividad por hora o si está afectada por cambios en la intensidad horaria.
\begin{table}[H]
\centering
\caption{Vaupés: PIB, ocupados, horas y productividad laboral, 2014--2024pr}
\label{tab:departamento_vaupes_productividad}
\scriptsize
\begin{tabular}{lrrr}
\toprule
Indicador & 2014 & 2024pr & Crec. anual \\
\midrule
PIB real & 0,2 & 0,3 & 2,24\% \\
Ocupados & 0,00 & 0,00 & 5,74\% \\
Horas semanales por trabajador & 45,6 & 43,0 & -0,57\% \\
PIB por trabajador & 81,9 & 58,5 & -3,31\% \\
PIB por hora & 34,6 & 26,2 & -2,75\% \\
\bottomrule
\end{tabular}
\end{table}

\subsection{Casanare}
\textbf{Entre 2014 y 2024, el PIB por hora trabajada de Casanare creció -3,33\% anual.} El PIB por trabajador creció -3,31\% anual, mientras que las horas semanales por trabajador cambiaron 0,01\% anual. Esta diferencia muestra si la productividad medida por trabajador se mueve en línea con la productividad por hora o si está afectada por cambios en la intensidad horaria.
\begin{table}[H]
\centering
\caption{Casanare: PIB, ocupados, horas y productividad laboral, 2014--2024pr}
\label{tab:departamento_casanare_productividad}
\scriptsize
\begin{tabular}{lrrr}
\toprule
Indicador & 2014 & 2024pr & Crec. anual \\
\midrule
PIB real & 13,4 & 13,1 & -0,22\% \\
Ocupados & 0,06 & 0,08 & 3,20\% \\
Horas semanales por trabajador & 46,8 & 46,9 & 0,01\% \\
PIB por trabajador & 231,4 & 165,2 & -3,31\% \\
PIB por hora & 95,1 & 67,8 & -3,33\% \\
\bottomrule
\end{tabular}
\end{table}

{\footnotesize Nota general: PIB real en billones de pesos constantes de 2015; ocupados en millones; PIB por trabajador en millones de pesos constantes de 2015; PIB por hora en miles de pesos constantes de 2015. Fuente: cálculos propios con DANE y GEIH.}